\subsection{Miscalibration Gap for Code Correctness}

Let $\{(p_i, r_i, y_i)\}_{i=1}^N$ denote a set of on-policy code patches, where
$r_i \in [0,10]$ is the model's self-rated correctness score for patch $p_i$ and
$y_i \in \{0,1\}$ indicates ground-truth correctness (1 if $p_i$ passes the SWE-bench tests).

We define the \textbf{miscalibration gap} as:
\begin{equation}
\mathrm{MiscalibrationGap}
\;=\;
\mathbb{E}[r_i \mid y_i = 0]
\;-\;
\mathbb{E}[r_i \mid y_i = 1].
\end{equation}

Intuitively, this quantity measures how well the model's scores separate failing patches from
passing patches. Values near $0$ (or positive) indicate poor separation: failing patches receive
scores comparable to (or higher than) passing patches.
We compute this quantity separately under each attribution condition (e.g., baseline vs.\ self-attributed)
and compare the resulting gaps to quantify how self-attribution affects calibration.
